\subsection{Finite covers couple the component and support fundamental groups}

\begin{statement}{C.2.UN.47}{Proved}{Positive geometric genus controls the whole support group}
Let $f:U\to\Delta$ be a proper holomorphic two-torus degeneration
with smooth connected total threefold, let $S=f^{-1}(0)_{\rm red}$,
and put $\Gamma=\pi_1(S)$. For an irreducible component $S_i$ let
$h_i:Z_i\to S_i\hookrightarrow U$ be a resolution of its normalization.
The following are simultaneous constraints on this actual fibre.

Every $Z_i$ with $p_g(Z_i)>0$ has $p_g(Z_i)=1$, and its actual
map induces a surjection
\[
                         \pi_1(Z_i)\longrightarrow\Gamma .
                                                        \tag{C.PF.1}
\]
Consequently every SNC dual complex of an embedded resolution of
the fibre is simply connected whenever such a component occurs.
A component with K3 minimal model forces $\Gamma=1$.
More generally,
\[
                    \chi(\OO_{Z_i})=2
                         \quad\Longrightarrow\quad\Gamma=1 .
                                                        \tag{C.PF.2}
\]
This includes the Kodaira-dimension-one components of holomorphic
Euler characteristic two in UN33.

For a component with $\chi(\OO_{Z_i})=1$, its fundamental-group
image in $\Gamma$ is zero, with just the following possible
group configuration:
\[
 \Gamma\simeq\Z/2\Z,\qquad
 \operatorname{im}\bigl(\pi_1(Z_i)\to\Gamma\bigr)=\Gamma .
                                                        \tag{C.PF.3}
\]
In that configuration the actual connected double cover of $Z_i$
induced from the tube has
$\chi(\OO)=2$, $p_g=1$ and $q=0$. At most one component with
$\chi(\OO)=1$ can have nontrivial fundamental-group image.
The assertion (C.PF.3) is a restriction on groups of actual fibres,
and is not an assertion that every surface or gluing with these
invariants occurs.
\end{statement}

\begin{proof}
\par\medskip\noindent\textbf{Every finite quotient gives another actual torus degeneration.}
Use the strong deformation retraction of UN44 on a sufficiently
small tube; its annular complement also retracts radially by proper
horizontal transport. Thus an open tube and its closed smaller
tube both have group $\Gamma$. Given a finite quotient
$\rho:\Gamma\twoheadrightarrow Q$ of order $N$, its connected
regular topological covering carries the unique complex structure
obtained by lifting the holomorphic coordinate charts. This gives
an actual finite \'{e}tale map
\[
                     p:U_Q\longrightarrow U,\qquad \deg p=N.
                                                        \tag{C.PF.4}
\]
It is proper: a finite covering of a locally compact Hausdorff
space is proper, as follows by covering a compact set with
finitely many relatively compact evenly covered neighbourhoods.
Consequently $f\circ p$ is proper. Its total space is smooth and
connected.

The punctured covering $U_Q^*$ is connected, since
$\pi_1(U^*)\to\Gamma\to Q$ is onto by UN44. Let $A$ be the image
of the nearby fibre group in $\Gamma$ and let $B=\rho(A)$.
The connected components over a fixed nearby fibre are indexed
by $Q/B$. The base generator permutes them transitively, and
$Q/B$ is cyclic, because $\Gamma/A$ is cyclic. If
$a=[Q:B]$, the connected base covering in the Stein factor of
$f\circ p$ is therefore the punctured degree-$a$ cover of the disc.
Its normal completion has one point above zero and is the disc
with actual coordinate map $t=s^a$. One can obtain this
identification by the classification of coverings of the
punctured disc and then extend the two inverse maps across zero:
each is bounded over a smaller disc, and the completed normal
one-dimensional spaces have the removable-singularity property.
Equivalently it is the normalization of this finite disc cover.
The proper Stein factor is this completion, because it is finite,
normal and agrees over the dense punctured disc. We have thus
constructed a proper map with connected fibres
\[
                        f_Q:U_Q\longrightarrow\Delta_s .
                                                        \tag{C.PF.5}
\]
Its reduced central support is exactly $(p^{-1}S)_{\rm red}$,
because the completed Stein disc has just one point over zero.
Every smooth general fibre of $f_Q$ is a connected finite \'{e}tale
cover of $F$. Writing $F=\C^2/\Lambda$, its actual lifted universal
cover identifies it with $\C^2/\Lambda'$ for a finite-index
sublattice $\Lambda'\subset\Lambda$. Hence (C.PF.5) is another
two-torus degeneration, to which UN28 applies. This comparison
does not require the original cover to be a pure base cover.

\par\medskip\noindent\textbf{The component pullbacks and their precise degree.}
Fix $Z=Z_i$ and take its fibre product with (C.PF.4). It is a
smooth compact, possibly disconnected surface
\[
               Z_Q=Z\times_U U_Q\longrightarrow Z
                       \quad\hbox{finite \'{e}tale of degree }N .
                                                        \tag{C.PF.6}
\]
Let $H_Q=\rho(h_{i*}\pi_1(Z))$, let $e=|H_Q|$, and put
$k=[Q:H_Q]=N/e$. Path lifting gives exactly $k$ connected
components $Z_1',\ldots,Z_k'$, each a connected cover of $Z$
of degree $e$.

These are resolutions of distinct actual normalization components
of the central support of $U_Q$. To verify this comparison, remove
the singular set of the original support and the exceptional
curves of $Z\to S_i$. Over the remaining dense open set the map
$Z\to U$ is its embedding as the smooth single-component locus.
Its base change is therefore an isomorphism onto the corresponding
open subset of the central support of $U_Q$. Removing those proper
analytic curves and points from a connected smooth $Z_j'$ leaves
it connected. Its proper image is consequently an irreducible
central component, and the displayed isomorphism over the dense
open set proves that $Z_j'$ maps birationally onto its normalization.
Different $Z_j'$ cannot give the same component, since their dense
open isomorphisms have disjoint images. Thus the compact-core
bound of UN28 gives the exact useful inequality
\[
                         \sum_{j=1}^k p_g(Z_j')\le1 .
                                                        \tag{C.PF.7}
\]
Other central components can only add nonnegative terms to that
bound.

Pullback of a holomorphic two-form under a finite \'{e}tale
cover is injective, by restriction to any open set on which
that cover is biholomorphic. If $p_g(Z)>0$, UN28 on the
original degeneration already gives $p_g(Z)=1$. Every
$Z_j'$ then contributes at least one in (C.PF.7), so $k=1$.
We have proved that
\[
       \rho(h_{i*}\pi_1(Z))=Q
                 \quad\hbox{for every finite quotient }\rho
                         \hbox{ of }\Gamma .           \tag{C.PF.8}
\]

\par\medskip\noindent\textbf{The group-theoretic passage retains the actual subgroup.}
UN44 supplies a normal finitely generated abelian subgroup
$A\subset\Gamma$ of finite index. Every subgroup $H$ of $\Gamma$
is separable by finite quotients in the sense needed here.
If $HA\ne\Gamma$, its image already fails to be all of
the finite quotient $\Gamma/A$. Otherwise, for a proper
$H$, choose $b\in A\setminus H$. The finitely generated
abelian group $A/(H\cap A)$ has a finite quotient detecting
the class of $b$; equivalently some positive integer $n$
satisfies
\[
                         b\notin(H\cap A)+nA .
                                                        \tag{C.PF.9}
\]
The subgroup $nA$ is invariant under every automorphism of $A$,
hence is normal in $\Gamma$, and $\Gamma/nA$ is finite.
Moreover $HN\cap A=(H\cap A)+nA$ for $N=nA$.
Thus the image of $H$ in this finite quotient misses $b$.
This proves directly that (C.PF.8) implies
$h_{i*}\pi_1(Z)=\Gamma$, establishing (C.PF.1).

The same elementary argument gives residual finiteness of
$\Gamma$. An element outside $A$ is detected in $\Gamma/A$.
A nonzero element of $A$ is detected in $A/nA$ for some $n$,
and therefore in the finite quotient $\Gamma/nA$.
No residual-finiteness assertion for a general degeneration
group is being imported.

\par\medskip\noindent\textbf{The dual complex and the K3 consequence.}
Choose any embedded resolution of the central divisor.
The strict transform of $S_i$ is a smooth resolution of its
normalization, with the same geometric genus as $Z$.
For completeness, the graph of the birational correspondence
between these two smooth resolutions has a common resolution
by point blowups. Pullback and Hartogs extension of holomorphic
two-form coefficients across a point show that each point blowup
preserves $p_g$. The same blowup maps induce isomorphisms on
$\pi_1$, by the $S^3$ and $\PP^1$ neighbourhood comparison.
The two maps from the common resolution to $U$ agree on a
dense open set and hence everywhere. Thus the image of the
strict-transform vertex group in $\Gamma$ is exactly the
image in (C.PF.1).

By (C.FG.14), the resolved support has group $\Gamma$.
By (C.FG.16), passage to its dual complex kills the image of
every component group. In particular it kills this surjective
vertex group, so the dual-complex fundamental group is trivial.
If $Z$ has K3 minimal model, its fundamental group is zero:
the K3 is simply connected and the point blowups do not change
its fundamental group. Equation (C.PF.1) then gives $\Gamma=1$.

\par\medskip\noindent\textbf{Holomorphic Euler characteristic gives stronger restrictions.}
The degree in (C.PF.6) also gives
\[
       \sum_{j=1}^k\chi(\OO_{Z_j'})=N\chi(\OO_Z).
                                                        \tag{C.PF.10}
\]
This follows, for example, by applying surface Riemann--Roch to
the actual \'{e}tale tangent-bundle pullback:
$12\chi(\OO)=\int(c_1^2+c_2)$, and integration of a pulled-back
top-degree class multiplies by the covering degree. Combining
this equality with
$p_g(Z_j')=\chi(\OO_{Z_j'})-1+q(Z_j')$ gives
\[
       1\ \ge\ \sum_j p_g(Z_j')
           \ \ge\ N\chi(\OO_Z)-k,\qquad k\le N .
                                                        \tag{C.PF.11}
\]
For $\chi(\OO_Z)=2$, the right side is at least $N$.
Every finite quotient of $\Gamma$ therefore has order one.
Residual finiteness proves (C.PF.2).

For $\chi(\OO_Z)=1$, a finite quotient in which $H_Q$ is
nontrivial has $e\ge2$ and satisfies
\[
                    1\ge N-\frac Ne\ge\frac N2 .
                                                        \tag{C.PF.12}
\]
It follows that $N=e=2$, so the only such finite quotient
is a quotient of order two. If the image
$H=h_{i*}\pi_1(Z)$ in $\Gamma$ is nontrivial, residual
finiteness provides one such quotient
$\rho:\Gamma\to\Z/2\Z$ detecting a nonidentity element of $H$.
If $\ker\rho$ contained a nonidentity element, another finite
quotient would detect that element. The image of the product
of these two quotient maps would then be a finite quotient
of order greater than two still detecting the original
element of $H$, contradicting (C.PF.12). Thus $\ker\rho=1$,
which proves precisely (C.PF.3).

In this case (C.PF.6) is a connected double cover of $Z$,
and (C.PF.10) gives $\chi=2$. Equation (C.PF.7) and
$p_g=\chi-1+q$ give $p_g=1$ and $q=0$. If two different
components with $\chi=1$ both had nontrivial image, the
same actual double cover of the tube would contain their
two distinct normalization resolutions, each of geometric
genus one. This contradicts UN28. It proves the final
joint assertion.
\end{proof}
